% Standalone problem document, generated from the adjacent README.md.
% Compile with XeLaTeX. Mathematical content is maintained in Markdown.
\documentclass[11pt,a4paper]{article}
\usepackage[fontsize=10.5pt]{scrextend}
\usepackage[margin=22mm,headheight=15pt,headsep=9mm,footskip=11mm]{geometry}
\usepackage{fontspec}
\setmainfont{texgyrepagella}[Extension=.otf,UprightFont=*-regular,BoldFont=*-bold,ItalicFont=*-italic,BoldItalicFont=*-bolditalic]
\setsansfont{texgyreheros}[Extension=.otf,UprightFont=*-regular,BoldFont=*-bold,ItalicFont=*-italic,BoldItalicFont=*-bolditalic]
\setmonofont{texgyrecursor}[Extension=.otf,UprightFont=*-regular,BoldFont=*-bold,ItalicFont=*-italic,BoldItalicFont=*-bolditalic,Scale=MatchLowercase]
\usepackage{amsmath,amssymb}
\usepackage{unicode-math}
\setmathfont{texgyrepagella-math.otf}
\usepackage{xcolor,microtype,enumitem,needspace,fancyhdr}
\usepackage{bookmark}
\definecolor{ink}{HTML}{17344B}
\definecolor{accent}{HTML}{197D80}
\definecolor{muted}{HTML}{596773}
\hypersetup{colorlinks=true,linkcolor=accent,urlcolor=accent,pdftitle={IE-13 - Sharp
growth for unequal lower and upper
bandwidths},pdfauthor={Open Problems in Numerical Linear Algebra}}
\urlstyle{same}
\setlength{\parindent}{0pt}
\setlength{\parskip}{5pt plus 1pt minus 1pt}
\linespread{1.04}
\setlength{\emergencystretch}{3em}
\widowpenalty=10000
\clubpenalty=10000
\displaywidowpenalty=10000
\allowdisplaybreaks[1]
\setlist{nosep,leftmargin=1.5em}
\providecommand{\tightlist}{\setlength{\itemsep}{0pt}\setlength{\parskip}{0pt}}
\providecommand{\pandocbounded}[1]{#1}
\pagestyle{fancy}
\fancyhf{}
\fancyhead[L]{\sffamily\footnotesize\color{muted}OPEN PROBLEMS IN NUMERICAL LINEAR ALGEBRA}
\fancyhead[R]{\sffamily\bfseries\footnotesize\color{accent}IE-13}
\fancyfoot[L]{\parbox[b]{0.9\textwidth}{\sffamily\fontsize{8}{10}\selectfont\color{muted}Editorial ratings; a bounded search cannot certify that a problem remains unsolved.\\Literature check: 2026-09-22}}
\fancyfoot[R]{\sffamily\footnotesize\color{muted}\thepage}
\renewcommand{\headrulewidth}{0.3pt}
\renewcommand{\headrule}{\hbox to\headwidth{\color{accent}\leaders\hrule height \headrulewidth\hfill}}
\makeatletter
\renewcommand{\section}{\Needspace{5\baselineskip}\@startsection{section}{1}{0pt}{12pt plus 2pt minus 2pt}{4pt}{\sffamily\bfseries\large\color{ink}}}
\renewcommand{\subsection}{\Needspace{4\baselineskip}\@startsection{subsection}{2}{0pt}{9pt plus 1pt minus 1pt}{3pt}{\sffamily\bfseries\normalsize\color{ink}}}
\makeatother
\setcounter{secnumdepth}{0}
\begin{document}
{\sffamily\small\color{muted}Linear systems and elimination\par}
\vspace{3pt}
{\raggedright\hyphenpenalty=10000\exhyphenpenalty=10000\sffamily\bfseries\fontsize{21}{24}\selectfont\color{ink}Sharp
growth for unequal lower and upper bandwidths\par}
\vspace{7pt}
{\sffamily\small\textbf{Difficulty:} challenging\quad\textbf{Importance:} interesting
to specialist\par}
{\sffamily\small\bfseries\color{accent}Status: Lean verified\par}
\vspace{4pt}
{\color{accent}\hrule height 0.8pt}
\vspace{6pt}
\subsection{Lean verification --- 22 September
2026}\label{lean-verification-22-september-2026}

The complete original unequal-bandwidth target is formally proved,
strengthened to every \(p,q\ge0\). The sharp growth is \(G(0,q)=1\) and
\(G(p,q)=h_{p+q}\) for \(p>0\), where \(h_t=0\) for \(t\le0\) and
\(h_t=1+\sum_{r=1}^p h_{t-r}\) otherwise. The universal upper bound
covers every nonsingular complex input with bandwidths at most \(p,q\)
in its original ordering, every \(n\ge1+\max(p,q)\), every permitted
GEPP tie path, and every intermediate entry. An explicit rational matrix
of order \(2p+q+1\) attains the bound along its actual full row-swap
path, including both zero-bandwidth endpoints. The entire growth set is
nonempty and bounded, has the stated greatest element, and has that
genuine real supremum.

All four exports passed LeanCert kernel audits, two independent
nonauthor final reviews, and
\href{https://github.com/sidneyholden1/OpenProblemsInNLA/actions/runs/35703222111}{isolated
Linux Comparator/default-kernel verification} with actual rejection and
sandbox controls.
\href{https://github.com/sidneyholden1/OpenProblemsInNLA/blob/main/linear-systems-and-elimination/IE-13/lean/README.md}{Proof
and reviews} ·
\href{https://github.com/sidneyholden1/OpenProblemsInNLA/tree/e8acacba4964ae8c103e808174cc6630d9bd00a1/linear-systems-and-elimination/IE-13/lean}{Immutable
proof revision} ·
\href{https://github.com/sidneyholden1/OpenProblemsInNLA/blob/main/docs/lean/verification-2026-09-22/IE-13/README.md}{Retained
evidence}.

Mathematics: Matthew J. Colbrook; Higham retains original-problem
credit. Formalization: Sidney Holden with OpenAI Codex assistance. No
normalization, assumed front invariant, real-only restriction,
preliminary reordering or deterministic tie rule is substituted. No
novelty, source-author endorsement or external human peer review is
claimed. The complete original target and earlier source record remain
below.

\subsection{Independently reviewed resolution -
2026-09-11}\label{independently-reviewed-resolution---2026-09-11}

\textbf{Sharp growth classification.} Theorem 1 and Sections 2-4 prove
\(G(0,q)=1\) and \(G(p,q)=h_{p+q}\) for \(p\ge1\), where \(h_t=0\) for
\(t\le0\) and \(h_t=1+\sum_{r=1}^p h_{t-r}\) otherwise. The upper bound
covers every complex input and admissible tie path, with growth over all
active entries in the fixed original ordering. A real nonsingular matrix
of order \(2p+q+1\) attains it, including zero upper bandwidth.

\textbf{Author:} Matthew J. Colbrook, Department of Applied Mathematics
and Theoretical Physics, University of Cambridge.
\href{https://github.com/ajt60gaibb/OpenProblemsInNLA/blob/main/references/colbrook-recovered-2026-09-11/manuscripts/IE-13.pdf}{Complete
manuscript},
\href{https://github.com/ajt60gaibb/OpenProblemsInNLA/blob/main/references/colbrook-recovered-2026-09-11/verification/reviews/IE-13-review.md}{independent
proof review}, and
\href{https://github.com/ajt60gaibb/OpenProblemsInNLA/blob/main/references/colbrook-recovered-2026-09-11/README.md}{submission
record}. The supplied notes were reconstructed with substantial AI
assistance. This is independent agent verification, not external human
peer review or formal proof-assistant certification; no novelty or
priority claim is made.

The difficulty, importance and rating rationale below are historical
assessments of the original open target. Original statements, references
and dated audits are preserved.

\textbf{Rating rationale:} Challenging reflects interacting fill-in and
pivot choices across arbitrary unequal bandwidths; specialist impact is
a sharp stability classification for banded elimination.

\newpage
\subsection{Problem statement}\label{problem-statement}

Work in exact arithmetic over \(\mathbb C\). For integers \(p,q\ge0\),
define

\[\mathcal B_n(p,q)=\{A\in\mathbb C^{n\times n}:\det A\ne0,\quad
a_{ij}=0\text{ if }i-j>p\text{ or }j-i>q\}.\]

Gaussian elimination with partial pivoting (GEPP) chooses a
largest-modulus entry in the active first column, moves its row to the
first position, and forms the trailing Schur complement. For a permitted
path \(\pi\), let \(S_1=A,\ldots,S_n\) be its active matrices and define

\[\rho(A,\pi)=\frac{\max_{1\le k\le n}\|S_k\|_{\max}}{\|A\|_{\max}},
\qquad \|M\|_{\max}=\max_{i,j}|m_{ij}|.\]

Determine, for all \(p\ne q\), the sharp bound independent of dimension,

\[G(p,q)=\sup_{n\ge1+\max(p,q)}\ \sup_{A\in\mathcal B_n(p,q)}\
\sup_{\pi\text{ permitted by GEPP}}\rho(A,\pi).\]

Identify the least universal upper bound and matching examples
approaching it. All tie choices are included. The bandwidth restrictions
apply in the given ordering; no reordering before GEPP is allowed.
Bandwidths are \emph{at most} \(p\) and \(q\); unequal pairs with one
zero bandwidth are included. For \(p\ge1\), the equal-bandwidth case is
known:

\[G(p,p)=2^{2p-1}-(p-1)2^{p-2}.\]

This asks how unequal bandwidths affect worst-case element growth in
banded linear solves. All unequal pairs constitute one problem.

\subsection{References}\label{references}

N. J. Higham,
\href{https://doi.org/10.1137/1.9780898718027}{\emph{Accuracy and
Stability of Numerical Algorithms}, second edition}, SIAM (2002),
Problem 9.15(a), p.~193; definitions and Theorem 9.11, pp.~172--173. Z.
Bohte, \href{https://doi.org/10.1093/imamat/16.2.133}{\emph{Bounds for
Rounding Errors in the Gaussian Elimination for Band Systems}}, J. Inst.
Maths. Applics. 16 (1975), 133--142.

\subsection{Earlier status check ---
2026-09-08}\label{earlier-status-check-2026-09-08}

The book's problem was checked directly. Its wording omits the field;
this entry adopts \(\mathbb C\) from its referenced Theorem 9.11.
Searches for \textit{Higham 9.15 growth},
\textit{growth factor lower bandwidth upper bandwidth}, and
\textit{Bohte Gaussian 1975} found no general unequal-bandwidth
solution. Shah--Urschel's
\href{https://arxiv.org/html/2608.19189v4}{August 31, 2026 revision},
Theorems 2.2--2.3, studies sparsity constraints without determining this
extremal function. No recent explicit reaffirmation of this particular
question's openness was located; the status rests on the original
question and this bounded later-literature check.

\subsection{Audit update --- 2026-09-10}\label{audit-update-2026-09-10}

Higham's
\href{https://pages.stat.wisc.edu/~bwu62/771/hingham2002.pdf}{Problem
9.15(a)}, p.~193, remains the explicit historical source.
Unequal-bandwidth and sharp-growth searches, including the
\href{https://arxiv.org/html/2608.19189v4}{August 2026 complete-pivoting
manuscript}, did not identify a sharp answer for this partial-pivoting
sparsity pattern. This is a bounded historical-source assessment,
without a recent explicit reaffirmation.
\end{document}
